\chapter{ Equivalencias }

  \section{Bicondicional}

  \begin{minipage}[t]{0.48\textwidth}
  \begin{definition}[title={Bicondicional}]
  $p \leftrightarrow q$ es verdad cuando $p$ y $q$ tienen el mismo valor
  de verdad en todas las interpretaciones, es decir:
  \begin{equation}
  p \leftrightarrow q \equiv (p \to q) \land (q \to p)
  \end{equation}
  \end{definition}
  \end{minipage}

  \begin{minipage}[t]{0.48\textwidth}
  \begin{tabular}{c|c|c|c|c}
  $p$ & $q$ & $p \to q$ & $q \to p$ & $(p \to q) \land (q \to p)$ \\
  \hline
  V & V & V & V & V \\
  F & V & V & F & F \\
  V & F & F & V & F \\
  F & F & V & V & V \\
  \end{tabular}
  \end{minipage}

  De forma que es verdadera cuando coinciden $p$ y $q$, así:

  \begin{equation}
  p \leftrightarrow q \equiv (p \land q) \lor (\sim p \land \sim q)
  \end{equation}

  Por ello, para probar su validez se suele probar $p \to q$ y posterior
  $q \to p$, pues los valores excluidos son $(V,F)$ y $(F,V)$ para cada
  implicación, de modo que $p \leftrightarrow q$ excluye ambos.
  Formalmente, probar que su caso de exclusión no puede ocurrir se deben
  mostrar ambos.

  De modo que los valores de exclusión son aquellos en los que el conector
  es falso por definición semántica.

  \subsection{Ejemplo: Triángulo Equilátero}

  \begin{proposition}[title={Ángulos de 60° implica equilátero}]
  Si el triángulo $ABC$ tiene 2 ángulos de $60°$ entonces es equilátero.
  \end{proposition}

  \begin{proof}[title={Construcción directa}]
  En un espacio euclídeo, todo triángulo tiene $180°$, de modo que el
  tercer ángulo es de $60°$ pues:
  \begin{align}
  x + 2(60) &= 180 \\
  x &= 60°
  \end{align}
  Luego, $ABC$ es equilátero. \qed
  \end{proof}

  \obs La demostración anterior muestra que tomando a $p$ como hipótesis,
  $q$ es el resultado por la restricción del modelo (el espacio euclídeo
  en este caso). De una forma indirecta se probó que $\sim(p \land \sim q)$
  es verdad, es decir, $p \to q$.

  Al probar que no puede pasar $p$ y $\sim q$ a la vez, se demuestra que
  se excluye el único caso de exclusión semántico. Por esto en una
  demostración, asumir $p$ y deducir $q$ es equivalente a probar que no
  puede ocurrir el caso $p \land \sim q$.

  \obs Si bien, en matemáticas no siempre las demostraciones usan estas
  reglas lógicas, hay equivalencias muy útiles.

  \section{Eliminación de la Implicación}

  \begin{definition}[title={Eliminación de la implicación}]
  \begin{equation}
  p \to q \equiv \sim p \lor q
  \end{equation}
  ya que su valor de exclusión es justo el mismo $(V, F)$, pues al
  evaluar, la disyunción ambos serían falsos.
  \end{definition}

  \section{Contraposición}

  \begin{minipage}[t]{0.48\textwidth}
  \begin{definition}[title={Contraposición}]
  Se parece mucho a equivalencia: si se asume que $q$ es $F$, entonces
  $p$ también debe serlo, pues la implicación excluye el caso en que $p$
  es $V$ y $q$ es $F$.
  \begin{equation}
  p \to q \equiv \sim q \to \sim p
  \end{equation}
  \end{definition}
  \end{minipage}

  \begin{minipage}[t]{0.48\textwidth}
  \begin{tabular}{c|c|c|c|c|c}
  $p$ & $q$ & $p \to q$ & $\sim p$ & $\sim q$ & $\sim q \to \sim p$ \\
  \hline
  V & V & V & F & F & V \\
  V & F & F & F & V & F \\
  F & V & V & V & F & V \\
  F & F & V & V & V & V \\
  \end{tabular}
  \end{minipage}

  Para este punto, el cómo se aborda una proposición de equivalencia es
  que en una demostración, se trata de probar que el caso de exclusión
  efectivamente no puede ocurrir, es decir, la proposición deberá ser una
  tautología para ser verdadera.

  \section{Leyes de De Morgan}

  \begin{theorem}[title={Leyes de De Morgan}]
  \begin{align}
  \sim(p \land q) &\equiv \sim p \lor \sim q \\
  \sim(p \lor q)  &\equiv \sim p \land \sim q
  \end{align}
  \end{theorem}

  Para que no ocurra $p \land q$, al menos una debe ser $F$. Ambas deben
  ser $F$ para que $p \lor q$ no se cumpla.

  \section{Formas Normales}

  Algo que puede ayudarnos a verificar factibilidad es convertir una
  proposición en su forma normal; esta es una forma estandarizada de
  escribir proposiciones usando solo ciertos conectores.

  \begin{definition}[title={Forma Normal Disyuntiva (FND)}]
  Es una conjunción de disyunciones:
  \begin{equation}
  (p_1 \lor q_1) \land (p_2 \lor q_2) \land \cdots
  \end{equation}
  \end{definition}

  \begin{definition}[title={Forma Normal Conjuntiva (FNC)}]
  Es una disyunción de conjunciones:
  \begin{equation}
  (p_1 \land q_1) \lor (p_2 \land q_2) \lor \cdots
  \end{equation}
  \end{definition}

  Toda proposición puede representarse por una FN. Para convertir una
  proposición a una forma normal, se siguen estos 3 pasos:

  \begin{enumerate}
  \item Se eliminan las implicaciones.
  \item Se expanden las negaciones usando las leyes de De Morgan, con la
  finalidad de poder agrupar linealmente.
  \item Se distribuyen los conectores.
  \end{enumerate}

  Por ejemplo:
  \begin{align}
  p \to (q \land r) &\equiv \sim p \lor (q \land r) \\
  \sim p \lor (q \land r) &\equiv (\sim p \lor q) \land (\sim p \lor r) \\
  \therefore\quad p \to (q \land r) &\equiv (\sim p \lor q) \land (\sim p \lor r)
  \end{align}